\section{Conclusion}
\label{sec:conclusion}

This paper asks whether PDUFA (1992) and GDUFA (2012) --- the two major user-fee-funded regulatory acceleration reforms at the FDA --- changed the controlled-substance share of original drug approvals. Using the complete Drugs@FDA database from 1939 through 2025 ($N = 25{,}908$) linked to DEA scheduling records, I estimate share-based interrupted time series, Poisson rate models with exposure, and stacked difference-in-differences specifications for both the NDA pathway (PDUFA) and the ANDA pathway (GDUFA), and supplement these with rate-versus-count diagnostics, group-specific linear time trends, and event-study DD designs.

The central finding is a precisely estimated null. In the most flexible specifications, the post-policy incidence-rate ratio for the controlled-substance share is $1.19$ for PDUFA on the NDA pathway (95\% CI $[0.59, 2.39]$) and $1.38$ for GDUFA on the ANDA pathway (95\% CI $[0.86, 2.21]$). The point estimates flip sign between less and more flexible specifications, which is itself diagnostic: when the empirical model is allowed to absorb the slow-moving secular movements in the CS and non-CS approval streams that have no plausible structural relationship to the user-fee policy variable, no level shift remains for that variable to explain. The apparent post-PDUFA NDA elevation in less flexible specifications is traceable to a separable late-2000s wave of branded extended-release and abuse-deterrent opioid formulations approved seventeen to twenty-one years after PDUFA enactment.

The findings carry two implications. First, the standard economic prediction --- that user-fee-funded review acceleration should disproportionately benefit drug classes with the largest expected revenue streams, including controlled substances --- does not find approval-side support in the data. The compositional variation that does exist is better explained by industry product cycles, post-Hatch--Waxman ANDA expansion, and the demand environment for prescription opioids than by review-time incentives at the regulatory margin. Second, this is not the same as saying that FDA regulatory policy is irrelevant to the controlled-substance supply chain. Approval is a necessary condition for market entry, and every product on the list in Table~\ref{tab:spike_drugs} required FDA clearance before it could reach prescribers. What the data do not support is the upstream causal chain from user-fee acceleration to approval composition.

For researchers and policymakers, the takeaway is that the causal pathway from regulatory acceleration to controlled-substance supply, if it exists, runs primarily through the innovation-investment channel --- firms deciding whether to develop CS products --- rather than through the review-time channel. Understanding that investment decision and how factors such as patent term, market exclusivity, and prescribing-market demand conditions shape it is likely to be more productive for policy design than further refinements to FDA review-time performance goals. Future work that links the upstream approval-side evidence reported here to downstream prescribing volumes and ultimately to health outcomes will provide a more complete picture of how pharmaceutical regulation shapes controlled-substance supply across the full market chain.
